\documentclass[12pt, a4paper]{article}
\usepackage{graphicx}
\usepackage{amsmath}
\usepackage{amssymb}
\graphicspath{ {./images/} }
\title{
  \begin{figure}[th]
    \centering
    \includegraphics[width=0.2\textwidth]{Univalle}
  \end{figure}
  \textbf{Universidad del Valle
    \\{\Large Facultad de ingeniería}
  \\{\large Ingeniería en sistemas}}}
\author{Cristian David Pacheco Torres
  \\ 2227437
  \\ Sebastian Molina Cuellar
  \\ xxxxx}
\date{Agosto 2022}
\begin{document}
\maketitle
Taller 1
\newpage{}
\begin{abstract}
Your abstract goes here functional programming
\end{abstract}
\newpage{}
\tableofcontents
\newpage{}
\section{Introduction}
A introduction
\section{Solucion ejercicios de programación funcional}
\subsection{Calcular el tamaño de una lista con un proceso iterativo}
\subsubsection{Proceso}
\subsubsection{Corrección}
\subsubsection{Casos de pruebas}
\subsection{Dividiendo una lista en dos sublistas a partir de un pivote}
\subsubsection{Proceso}
\subsubsection{Corrección}

Sea $$f: List[Int] \rightarrow List[Int]$$ una funcion con dominio en el conjunto S finito de k >= 1 elementos, con codominio en el conjunto M de listas finitas
tal que m contenido en S, y por tanto l <=
Sea L = List(a_i, ..., a_k) un elemento de S
Se definine entonces las caracterizacion del proceso iterativo
  s = (L, M, v) donde M y L son listas y v es el valor pivote con el cual se compara
  s_o = (L, [], v)  L la lista original, v el valor pivote
  s_i = s_f si en s = (l, m, v) l es la lista vacia
  inv(s_i) = \equiv l = List(a_i, a_{i+1}...,a_k) ^ m \in List(a_1, ...., a_i-1) ^ m = f(List(a_1, ...., a_i-1))
  T(s_i) = (l(a_{i+1},...,k)), m, v)  donde si a_i < v \tehn a_i::m | m \in L(a_1,...,a_{i-1}) sino m \in L(a_1,...,a_{i-1})

Sea $$k: List[Int] \rightarrow Int$$, un numero entero k | 0 <= k <= tamI(l), un numero entero que representa el elemento en la posicion k de la lista


s = ()

\subsubsection{Casos de pruebas}
\subsection{Calculando el k-ésimo elemento de una lista}
\subsubsection{Proceso}
\subsubsection{Corrección}
\subsubsection{Casos de pruebas}
\section{Conclusion}
La conclusion
\newpage{}
\begin{displaymath}
 a = \sum F \dot m = \frac{dv}{dt}
\end{displaymath}

\end{document}
